\heading{高等数学A（II）-23春-期末}
\noteinfo{题目来源于赛艇，答案由AI生成。}

\begin{question}[三个初值问题]

\begin{enumerate}[label=(\arabic*),leftmargin=2.6em,itemsep=0.25em,topsep=0.2em]
\item 求微分方程 $y'-y=e^x,\ y(0)=0$ 的解。
\item 求微分方程 $y'=y^{1/3},\ y(0)=0$ 的解。
\item 求微分方程 $y''-2y'+y=e^{\alpha x},\ y(0)=0,\ y'(0)=0$ 的解，其中 $\alpha$ 为常数。
\end{enumerate}

\begin{solution}

\begin{enumerate}[label=(\arabic*),leftmargin=2.6em,itemsep=0.25em,topsep=0.2em]
\item 乘积分因子 $e^{-x}$：
\[
 (e^{-x}y)'=1
 \Longrightarrow e^{-x}y=x+C.
\]
由 $y(0)=0$ 得 $C=0$，故
\ansmath{y=xe^x.}

\item 方程可分离：
\[
 y^{-1/3}dy=dx
 \Longrightarrow \frac32 y^{2/3}=x+C.
\]
由于右端在 $y=0$ 处不满足 Lipschitz 条件，初值问题\emph{不唯一}。典型解包括
\[
 y\equiv 0,
\]
以及对任意 $c\ge 0$，
\[
 y(x)=
 \begin{cases}
 0,&x\le c,\\
 \bigl(\tfrac23(x-c)\bigr)^{3/2},&x\ge c.
 \end{cases}
\]

\item 若 $\alpha\ne 1$，取特解 $y_p=\dfrac{e^{\alpha x}}{(\alpha-1)^2}$，再加齐次解 $(C_1+C_2x)e^x$。由初值得
\ansmath{y(x)=\frac{e^{\alpha x}-e^x-(\alpha-1)xe^x}{(\alpha-1)^2},\qquad \alpha\ne1.}
若 $\alpha=1$，为二重共振，取特解 $y_p=\dfrac12 x^2e^x$，且自动满足初值，所以
\ansmath{y(x)=\frac12 x^2e^x,\qquad \alpha=1.}
\end{enumerate}

\end{solution}
\end{question}

\begin{question}[一致收敛]

讨论函数项级数
\[
\sum_{n=1}^{\infty}\frac{\sin nx}{n^p},\qquad p>0
\]
在区间 $[1,2]$ 内是否一致收敛。

\begin{solution}

利用 Dirichlet 判别法。对任意 $x\in[1,2]$，部分和
\[
\sum_{k=1}^{N}\sin kx
\]
可写成
\[
\sum_{k=1}^{N}\sin kx=
\frac{\sin\frac{Nx}{2}\,\sin\frac{(N+1)x}{2}}{\sin\frac{x}{2}}.
\]
由于在区间 $[1,2]$ 上有 $\sin(x/2)$ 远离 $0$，故上述部分和关于 $x$ 一致有界。另一方面，$1/n^p$ 单调趋于 $0$。因此对任意 $p>0$，级数在 $[1,2]$ 上一致收敛。

所以答案是
\ansmath{\text{对一切 }p>0,\ \sum_{n=1}^{\infty}\frac{\sin nx}{n^p}\text{ 在 }[1,2]\text{上一致收敛。}}

\end{solution}
\end{question}

\begin{question}[Fourier 级数]

求 $f(x)=\sin^2x$ 在 $[-\pi,\pi]$ 上的 Fourier 级数。

\begin{solution}

利用恒等式
\[
\sin^2x=\frac{1-\cos 2x}{2},
\]
因此 Fourier 级数就是它本身：
\ansmath{\sin^2x=\frac12-\frac12\cos 2x.}

\end{solution}
\end{question}

\begin{question}[参数型级数的收敛域]

讨论函数项级数
\[
\sum_{n=1}^{\infty}\frac{(-1)^n}{n^x+n}
\]
的收敛域。

\begin{solution}

对任意实数 $x$，都有 $n^x+n\to\infty$，故通项趋于 $0$。并且
\[
 a_n(x)=\frac{1}{n^x+n}>0
\]
对固定 $x$ 而言最终单调下降，因此可用 Leibniz 判别法知：该级数对\emph{所有}实数 $x$ 都收敛。

再讨论绝对收敛：
\[
\sum_{n=1}^{\infty}\left|\frac{(-1)^n}{n^x+n}\right|
=\sum_{n=1}^{\infty}\frac{1}{n^x+n}.
\]
\begin{enumerate}[label=(\arabic*),leftmargin=2.6em,itemsep=0.25em,topsep=0.2em]
\item 当 $x>1$ 时，分母主项为 $n^x$，故级数与 $\sum n^{-x}$ 同敛散，绝对收敛。
\item 当 $x\le 1$ 时，分母主项为 $n$，故与调和级数同阶，不绝对收敛。
\end{enumerate}
故结论为
\ansmath{\text{收敛域为 }\mathbb R;\text{其中 }x>1\text{ 绝对收敛，}x\le1\text{ 条件收敛。}}

\end{solution}
\end{question}

\begin{question}[幂级数展开]

设
\[
 f(x)=\ln(1+x+x^2+x^3+x^4).
\]
\begin{enumerate}[label=(\arabic*),leftmargin=2.6em,itemsep=0.25em,topsep=0.2em]
\item 求 $f(x)$ 在 $x_0=0$ 处的幂级数展开 $\sum_{n=1}^{\infty}a_nx^n$。
\item 确定该幂级数的收敛域。
\item 计算 $\sum_{n=1}^{\infty}a_n$ 的值。
\end{enumerate}

\begin{solution}

注意
\[
1+x+x^2+x^3+x^4=\frac{1-x^5}{1-x}\qquad (x\ne1).
\]
因此
\[
 f(x)=\ln(1-x^5)-\ln(1-x).
\]
在 $|x|<1$ 内展开：
\[
 f(x)=-\sum_{m=1}^{\infty}\frac{x^{5m}}{m}+\sum_{m=1}^{\infty}\frac{x^m}{m}.
\]
所以
\[
 a_n=
 \begin{cases}
 \dfrac1n, & 5\nmid n,\\
 -\dfrac4n, & 5\mid n.
 \end{cases}
\]
这给出了幂级数展开。

因为奇点是五次单位根（除 $1$ 外）与 $x=1$，都位于单位圆上，所以收敛半径为 $1$。在实轴上，$x=1$ 与 $x=-1$ 处也收敛，因此可写为
\ansmath{\text{实收敛域为 }[-1,1].}

最后
\[
\sum_{n=1}^{\infty}a_n=f(1^-)=\ln(1+1+1+1+1)=\ln 5.
\]
故
\ansmath{\sum_{n=1}^{\infty}a_n=\ln 5.}

\end{solution}
\end{question}

\begin{question}[积分与极限]

\begin{enumerate}[label=(\arabic*),leftmargin=2.6em,itemsep=0.25em,topsep=0.2em]
\item 设 $0\le a<b$，计算含参变量积分
\[
\int_a^b\frac{x}{y^2+x^2}\,dy,\qquad x>0.
\]
\item 计算极限
\[
\lim_{x\to+\infty}\sum_{n=1}^{+\infty}\frac{x}{n^2+x^2}.
\]
\end{enumerate}

\begin{solution}

\begin{enumerate}[label=(\arabic*),leftmargin=2.6em,itemsep=0.25em,topsep=0.2em]
\item 直接积分：
\[
\int\frac{x}{y^2+x^2}\,dy=\arctan\frac{y}{x}.
\]
因此
\ansmath{\int_a^b\frac{x}{y^2+x^2}\,dy=\arctan\frac{b}{x}-\arctan\frac{a}{x}.}

\item 写成 Riemann 和：
\[
\sum_{n=1}^{\infty}\frac{x}{n^2+x^2}
=\frac1x\sum_{n=1}^{\infty}\frac{1}{1+(n/x)^2}.
\]
当 $x\to\infty$ 时，这趋于
\[
\int_0^{\infty}\frac{dt}{1+t^2}=\frac{\pi}{2}.
\]
故
\ansmath{\lim_{x\to+\infty}\sum_{n=1}^{+\infty}\frac{x}{n^2+x^2}=\frac{\pi}{2}.}
\end{enumerate}

\end{solution}
\end{question}

\begin{question}[级数敛散性]

设 $a\ne0$、$p>0$，讨论级数
\[
\sum_{n=2}^{+\infty}\frac{\sin\bigl(\pi\sqrt{n^2+a^2}\bigr)}{(\ln n)^p}
\]
的敛散性;若收敛，问是绝对收敛还是条件收敛？

\begin{solution}

当 $n\to\infty$ 时，
\[
\sqrt{n^2+a^2}=n+\frac{a^2}{2n}+O(n^{-3}).
\]
因此
\[
\sin\bigl(\pi\sqrt{n^2+a^2}\bigr)
=\sin\left(\pi n+\frac{\pi a^2}{2n}+O(n^{-3})\right)
=(-1)^n\left(\frac{\pi a^2}{2n}+O(n^{-3})\right).
\]
所以原级数通项与
\[
(-1)^n\frac{C}{n(\ln n)^p}
\]
同阶。于是：
\begin{enumerate}[label=(\arabic*),leftmargin=2.6em,itemsep=0.25em,topsep=0.2em]
\item 绝对值级数与
\[
\sum\frac{1}{n(\ln n)^p}
\]
同敛散，因此当且仅当 $p>1$ 时绝对收敛;
\item 当 $0<p\le1$ 时，绝对值级数发散，但原级数为交错型且振幅单调趋于 $0$，故条件收敛。
\end{enumerate}
结论：
\ansmath{p>1\text{ 时绝对收敛; }0<p\le1\text{ 时条件收敛。}}

\end{solution}
\end{question}

\begin{question}[特殊初值问题]

设 $y=y(x)$ 是如下初值问题的解：
\[
 y'=\sqrt{1-y^2},\qquad y(0)=1.
\]
试计算 $y'(1)$ 的值。

\begin{solution}

因为右端定义域要求 $|y|\le1$，且
\[
 y'(x)=\sqrt{1-y(x)^2}\ge0,
\]
所以解不可能从 $y(0)=1$ 再增大;唯一可能是恒有 $y\equiv1$。于是
\ansmath{y'(1)=0.}

\end{solution}
\end{question}

\begin{question}[可分离变量微分方程]

求初值问题
\[
 y'=\frac{\sqrt{1+y^2}}{\sqrt{1+x^2}},\qquad y(1)=1
\]
的解。

\begin{solution}

分离变量：
\[
\frac{dy}{\sqrt{1+y^2}}=\frac{dx}{\sqrt{1+x^2}}.
\]
积分得
\[
\operatorname{arsinh}y=\operatorname{arsinh}x+C.
\]
由 $y(1)=1$ 得 $C=0$，所以
\ansmath{\operatorname{arsinh}y=\operatorname{arsinh}x\Longrightarrow y=x.}

\end{solution}
\end{question}

\clearpage

